\documentclass[11pt,letterpaper]{article}
\usepackage[T1]{fontenc}
\usepackage[utf8]{inputenc}
\usepackage{lmodern}
\usepackage[margin=1in,headheight=15pt]{geometry}
\usepackage{amsmath,amssymb,amsthm,mathtools,mathrsfs}
\usepackage{microtype,booktabs,longtable,array,enumitem}
\usepackage[dvipsnames]{xcolor}
\usepackage{fancyhdr}
\usepackage[colorlinks=true,linkcolor=MidnightBlue,citecolor=MidnightBlue,urlcolor=MidnightBlue]{hyperref}
\usepackage[nameinlink,noabbrev]{cleveref}
\hypersetup{pdftitle={Finite Intersections in Surcomplex Analysis},pdfauthor={Research manuscript},pdfsubject={Hahn-coherent division, residue duality, and sharp root stability}}
\pagestyle{fancy}\fancyhf{}
\fancyhead[L]{\small Finite intersections in surcomplex analysis}
\fancyhead[R]{\small Division, duality, stability}
\fancyfoot[C]{\thepage}
\setlength{\parskip}{0.3em}
\setlength{\parindent}{1.25em}
\setlist{nosep,leftmargin=2em}
\newtheorem{theorem}{Theorem}[section]
\newtheorem{lemma}[theorem]{Lemma}
\newtheorem{proposition}[theorem]{Proposition}
\newtheorem{corollary}[theorem]{Corollary}
\theoremstyle{definition}
\newtheorem{definition}[theorem]{Definition}
\newtheorem{example}[theorem]{Example}
\theoremstyle{remark}
\newtheorem{remark}[theorem]{Remark}
\newcommand{\No}{\mathbf{No}}
\newcommand{\SC}{\mathbf{SC}}
\newcommand{\C}{\mathbb C}
\newcommand{\N}{\mathbb N}
\newcommand{\Z}{\mathbb Z}
\newcommand{\Q}{\mathbb Q}
\newcommand{\OO}{\mathcal O}
\newcommand{\HH}{\mathcal H}
\newcommand{\mm}{\mathfrak m}
\newcommand{\BB}{\mathcal B}
\newcommand{\PP}{\mathcal P}
\newcommand{\RR}{\mathscr R}
\newcommand{\sh}{^{\sharp}}
\newcommand{\supp}{\operatorname{supp}}
\newcommand{\st}{\operatorname{st}}
\newcommand{\NF}{\operatorname{NF}}
\newcommand{\Tr}{\operatorname{Tr}}
\newcommand{\Span}{\operatorname{span}}
\newcommand{\id}{\operatorname{id}}
\newcommand{\Res}{\operatorname{Res}}
\newcommand{\val}{\operatorname{v}}
\newcommand{\ord}{\operatorname{ord}}
\newcommand{\HInt}{\mathop{\int^{\!H}}\nolimits}
\newcommand{\dz}{\,dz_1\wedge\cdots\wedge dz_n}
\newcolumntype{L}[1]{>{\raggedright\arraybackslash}p{#1}}
\numberwithin{equation}{section}
\begin{document}
\begin{titlepage}
\centering
\vspace*{0.7cm}
{\Large\scshape A research extension of Hahn-coherent analysis\par}
\vspace{0.8cm}
{\Huge\bfseries Finite Intersections\par in Surcomplex Analysis\par}
\vspace{0.65cm}
{\LARGE Support-Controlled Division,\par Residue Duality, and Sharp Root Stability\par}
\vspace{0.8cm}
{\large September 21, 2026\par}
\vspace{0.7cm}
\begin{minipage}{0.92\textwidth}
\begin{abstract}
We develop a several-variable finite-intersection theory for holomorphic functions with well-ordered Hahn support, evaluated at surcomplex infinitesimal points. For a coupled system
\[
 F_i(s,z)=z_i^{d_i}+E_i(s,z),\qquad v_H(E_i)>0,
\]
we construct an explicit division operator on a fixed ordinary polydisk. The quotient is free of rank $N=\prod_i d_i$ over the Hahn-holomorphic parameter algebra, with the usual box monomials as a basis. An explicit contraction of the entire perturbed Koszul complex supplies both ideal-membership certificates and exact support bounds.

We prove that this algebra is the full intersection algebra of the actual surcomplex zero cluster: its local factors have the formal analytic intersection multiplicities, whose sum is exactly $N$. The associated map on an infinitesimal polydisk is surjective and finite of constant multiplicity. A coefficientwise multidimensional contour functional yields a perfect residue pairing, a B\'ezoutian reproducing kernel, the Jacobian trace formula, and residues at infinitesimally separated points.

A second theorem gives sharp valuation-scale stability. If a simple root has inverse-Jacobian loss $\kappa$, a perturbation of valuation $\sigma>2\kappa$ has a unique nearby simple root, with explicit first-order and remainder bounds; the strict threshold cannot be weakened uniformly. All infinite constructions use finite contributions at each exponent, not sequential convergence. A coupled two-variable family illustrates flat collision, discriminants, and cancellation of infinite residues.
\end{abstract}
\end{minipage}
\vfill
\begin{minipage}{0.92\textwidth}\small
\textbf{Status.} The precise support-controlled package is proposed as a new extension of the supplied manuscript. Its classical ingredients are identified explicitly. The literature search did not establish priority, and this is not a claim to have settled a named published conjecture. The proofs are mathematical arguments, supplemented by exact symbolic checks, not a machine-verified development.
\end{minipage}
\end{titlepage}
\tableofcontents
\newpage

\section{The research problem and the results}\label{sec:intro}
The supplied manuscript \cite{Supplied} established a one-variable theory of Hahn-coherent holomorphy, including preparation, root lifting, and moving-pole residues. It explicitly left several-variable coherent division and finite mappings as further targets. This article treats a substantial, precisely delimited part of that problem: arbitrary positive-Hahn-support perturbations of a coordinate-monomial complete intersection, including unrestricted coupling among the variables and ordinary holomorphic parameters.

The starting system $z_i^{d_i}=0$ is elementary, but its perturbations need not be polynomial, finitely supported, or controlled by a single infinitesimal parameter. Their coefficient functions may grow arbitrarily fast as the Hahn exponent varies. The point is to construct one analytic quotient and one residue theory that remain valid for all these perturbations on the original ordinary domain. An iteration that continually shrinks domains, or that assumes a real-valued valuation, does not suffice.

\subsection{Three principal statements}
Let $\Gamma\subseteq\No$ be a set-sized divisible ordered additive subgroup and put
\[
 K=\C((t^\Gamma)),\qquad t^\gamma=\omega^{-\gamma}.
\]
This is an algebraically closed subfield of $\SC=\No[i]$. Let $D\subset\C^n$ be a polydisk centered at zero, $V$ an ordinary parameter domain, and
\[
 \HH_\Gamma(V\times D)=\OO(V\times D)((t^\Gamma)).
\]
The notation means a well-ordered Hahn family of holomorphic coefficient functions, not a completed tensor product chosen implicitly. The definitions are made precise in \cref{sec:setting}.

\paragraph{Finite free division.}
For $F_i=z_i^{d_i}+E_i$, with every nonzero $E_i$ supported at strictly positive exponents, there is a canonical remainder
\[
 \NF_F(H)=\sum_{0\le\alpha_i<d_i}c_\alpha(s)z^\alpha,
 \qquad c_\alpha\in\HH_\Gamma(V),
\]
and a constructed division identity $H=\NF_F(H)+\sum_i F_iG_i$ on the same $V\times D$. The remainder is unique, even though a tuple of quotients generally is not. If $S=\supp H$ and $T$ is the union of the perturbation supports, all constructed coefficients have support in $S+T^*$. This gives a free quotient of rank $N=\prod_i d_i$ and an explicit contraction in every positive Koszul degree.

\paragraph{Intersection algebra and residue duality.}
With parameters fixed, all actual zeros lie in the infinitesimal polydisk $\mm_K^n$. Their formal local intersection algebras have dimensions $\mu_a$ satisfying
\[
 \sum_{F(a)=0}\mu_a=N.
\]
A Hahn contour functional $\RR_F$ makes the quotient a Frobenius algebra: $(u,v)\mapsto\RR_F(uv)$ is nondegenerate. If $M_H$ denotes multiplication by $H$ and $J_F$ the Jacobian determinant, then
\begin{equation}\label{eq:introtrace}
 \Tr(M_H)=\RR_F(HJ_F)
          =\sum_{F(a)=0}\mu_a H\sh(a).
\end{equation}
At simple zeros the localized residue is $H\sh(a)/J_F\sh(a)$, including when this value is infinite.

\paragraph{Sharp root stability.}
For a simple zero $a$, put $A=DF\sh(a)$ and $\kappa=\max\{0,-v(A^{-1})\}$. If $\widetilde F=F+\Delta$ and $v_H(\Delta)\ge\sigma>2\kappa$, there is exactly one zero $a+h$ in the ball $v(h)>\kappa$, and
\begin{equation}\label{eq:introstability}
 v(h)\ge\sigma-\kappa,\qquad
 v\bigl(h+A^{-1}\Delta\sh(a)\bigr)\ge2\sigma-3\kappa.
\end{equation}
The new zero is simple. A quadratic example shows that equality $\sigma=2\kappa$ can merge two simple roots, so strictness is necessary for a uniform statement.

\subsection{What is and is not being claimed as new}
Hahn summability, local surreal analytic evaluation, homological perturbation, finite-dimensional intersection algebras, Grothendieck residues, B\'ezoutians, and Hensel-type conditioning are established subjects. In particular, the conceptual relation between normal forms, residues, and traces is already developed by Cattani--Dickenstein--Sturmfels \cite{CDS}; it is not an original discovery of this article. The proposed contribution is their \emph{support-controlled analytic realization and compatibility} in the category specified here, together with the fixed-domain parameter theorem and the explicit arbitrary-rank stability statement.

The result is stronger than extending a finite polynomial identity by a change of coefficient field. The inputs can be infinite Hahn families of nonpolynomial holomorphic functions. One must prove that division preserves their common ordinary domain, that abstract algebra characters coincide with actual surcomplex evaluation, and that residue identities survive summation without assuming rank-one convergence. These issues receive separate proofs.

We do not assert a theorem for every isolated ordinary complete intersection, a global transcendental theory on all of $\SC^n$, or a resolution of all the final questions in \cite{Supplied}. The coordinate-monomial reduction is an explicit hypothesis, not a concealed generality assumption. A precise extension criterion and remaining questions appear in \cref{sec:scope}.

\subsection{Relation to published work}
Alling's analysis and the restricted analytic structure of surreal fields supply the background for evaluation \cite{Alling,vdDE}. Poonen's algebraic-closure theorem supplies the algebraically closed Hahn workspace \cite{Poonen}. The conjugation argument below is a particularly transparent specialization of homological perturbation, a general method discussed by Crainic \cite{Crainic}. The multiplication-matrix viewpoint has a classical account in Cox \cite{Cox}.

The terminology needs care. M\"uller--Strohmaier study Hahn-holomorphic and Hahn-meromorphic expansions of ordinary complex functions, using normal convergence and generalized monomials on a logarithmic covering \cite{MS}. Our coefficient functions live on an ordinary analytic base, but their Hahn weights are fixed surreal monomials and their evaluations are surcomplex. The two settings are related in spirit, not identical by definition. Berarducci--Mantova and Costin--Ehrlich address other, important surreal function and integration frameworks \cite{BM,CE}. We do not claim a comparison of all these constructions.

\section{Hahn-coherent analytic data}\label{sec:setting}
\subsection{Fields, valuation, and infinitesimal evaluation}
A series in $K$ is an expression
\[
 c=\sum_{\gamma\in S}c_\gamma t^\gamma,
 \quad c_\gamma\in\C,\quad S\subseteq\Gamma\text{ well ordered}.
\]
Set $v(c)=\min\supp(c)$ when $c\ne0$, and $v(0)=+\infty$. Then
\[
 v(cd)=v(c)+v(d),\qquad v(c+d)\ge\min\{v(c),v(d)\}.
\]
Write $K^\circ=\{c:v(c)\ge0\}$ and $\mm_K=\{c:v(c)>0\}$. The residue field is $\C$. Under the normal-form embedding in $\SC$, $K^\circ$ consists of finite elements and $\mm_K$ of infinitesimals. Divisibility of $\Gamma$ and algebraic closedness of $\C$ imply algebraic closedness of $K$ \cite[Corollary 4]{Poonen}.

For a connected ordinary domain $U\subseteq\C^m$, define $\HH_\Gamma(U)$ to consist of
\begin{equation}\label{eq:coherent}
 H=\sum_{\gamma\in S}h_\gamma t^\gamma,
 \qquad h_\gamma\in\OO(U),
\end{equation}
with common well-ordered support $S$. Zero coefficient functions are omitted. Its valuation $v_H$ is the least exponent of a nonzero coefficient function. A point $c+\varepsilon$, where $c\in U$ is ordinary and $\varepsilon\in\mm_K^m$, has evaluation
\begin{equation}\label{eq:evaluation}
 H\sh(c+\varepsilon)=
 \sum_{\gamma,\alpha\in\N^m}
 t^\gamma\frac{\partial^\alpha h_\gamma(c)}{\alpha!}\varepsilon^\alpha.
\end{equation}
The thickened domain is denoted $U_K\sh$. Evaluation is allowed in a larger divisible Hahn workspace as well. All index sets in this article are sets, not proper classes.

\begin{definition}[Strong summability]
A family of Hahn series is strongly summable if the union of its supports is well ordered and only finitely many family members contribute at any one exponent. Its sum is formed by finite coefficient addition at each exponent. The same definition applies to Hahn families with holomorphic coefficient functions, vectors, matrices, and finite exterior powers.
\end{definition}

\begin{lemma}[Support lemma]\label{lem:support}
Let $S,T$ be well-ordered subsets of an ordered abelian group. Then $S+T$ is well ordered and each exponent has finitely many representations as $s+t$. If $T$ consists of strictly positive elements, its additive monoid
\[
 T^*=\{\tau_1+\cdots+\tau_k:k\ge0,\ \tau_j\in T\}
\]
is well ordered and each exponent has finitely many representations as a finite word in $T$, up to permutation. Consequently each exponent has finitely many representations as an element of $S$ plus a finite word in $T$.
\end{lemma}
\begin{proof}
The first assertion follows from the impossibility of a decreasing subsequence in a well-ordered set: an infinite descending sequence of sums, or infinitely many pairs with the same sum, would force one coordinate to decrease along an infinite subsequence. The monoid assertion is the Hahn--Neumann lemma \cite{Neumann}. One proof uses the well-quasi-ordering of finite words in a well-ordered positive alphabet under coordinatewise embedding. An embedding cannot decrease the sum; equality forces equal words. This excludes descending sums and infinitely many words of the same sum. There are only finitely many permutations of each finite word. Apply the first assertion to $S$ and $T^*$ for the last statement.
\end{proof}

This proves summability of \eqref{eq:evaluation}: its support lies in $S+W^*$, where $W$ is the finite union of the supports of the infinitesimal coordinates. Finite contributions also justify multiplication and the usual Taylor identities. Evaluating at all ordinary points proves injectivity by uniqueness of normal forms. Differentiation acts on the coefficient functions, leaving the monomials $t^\gamma$ constant.

\subsection{Why support, rather than sequential convergence, is essential}
If $0<\alpha\ll\beta$ in $\Gamma$, then $k\alpha<\beta$ for every ordinary integer $k$. Thus the remainder valuations $k\alpha$ of a geometric-series calculation need not tend cofinally to $+\infty$. In the full surreal fine topology the obstruction is stronger: a set of nonzero positive distances has a smaller positive surreal lower bound, so any convergent set-indexed net is eventually its limit. The full proof and the resulting distinctions between coherence and fine analyticity are in \cite{Supplied}.

We will never prove an infinite construction by writing ``each step increases valuation, hence the iteration converges.'' Instead we exhibit a support $S+T^*$ and prove finite contribution at every exponent. The following operator version is the workhorse.

\begin{lemma}[Support-controlled operator inversion]\label{lem:operator}
Suppose a linear operator $L$ on a finite direct sum of Hahn-holomorphic spaces is built from finitely many coefficientwise complex-linear operators and multiplication by positive-support Hahn data with support union $T$. Assume each application of $L$ introduces at least one such multiplier. Then
\[
 (1+L)^{-1}H=\sum_{k\ge0}(-L)^kH
\]
is strongly summable. Its support is contained in $\supp(H)+T^*$. It is a two-sided inverse, provided the coefficientwise operators are defined on the common ordinary domains.
\end{lemma}
\begin{proof}
Expand each $L^k$ into words of coefficientwise operators and positive-support multipliers. The number of operator patterns for a fixed finite length is finite. At a fixed exponent, \cref{lem:support} permits only finitely many choices of the initial exponent, multiplier exponents, and word lengths. The resulting coefficient is therefore a finite composition of operators applied to holomorphic functions. Telescoping the geometric identity at that coefficient proves both inverse identities. No limit of the partial sums is used.
\end{proof}

\section{A fixed-domain contraction for a monomial intersection}\label{sec:contraction}
Let
\[
 D=\prod_{i=1}^n\{z_i\in\C:|z_i|<r_i\},\qquad r_i>0,
\]
and fix positive integers $d_1,\ldots,d_n$. Let $V\subseteq\C^p$ be a connected ordinary domain; the case without parameters means $\OO(V)=\C$. Set
\[
 R=\OO(V\times D),\quad B=\OO(V),\quad
 \BB=\{\alpha\in\N^n:0\le\alpha_i<d_i\},\quad N=|\BB|.
\]
Let $\PP=\bigoplus_{\alpha\in\BB}Bz^\alpha$.

\subsection{Taylor division operators}
For $f\in R$, define
\begin{align}
 R_i f&=\sum_{j=0}^{d_i-1}\frac{z_i^j}{j!}
   (\partial_i^j f)(s,z_1,\ldots,0,\ldots,z_n),\label{eq:Ri}\\
 D_i f&=\frac{f-R_i f}{z_i^{d_i}}.\label{eq:Di}
\end{align}
The apparent singularity of $D_i f$ is removable, so both operators preserve holomorphy on the \emph{same} $V\times D$. They are $B$-linear. For distinct indices they commute; for a fixed index,
\begin{equation}\label{eq:divrelations}
 z_i^{d_i}D_i=1-R_i,\qquad
 D_i z_i^{d_i}=1,\qquad D_iR_i=0.
\end{equation}
The box projection is $p=R_1\cdots R_n:R\to\PP$.

\subsection{The Koszul complex with an explicit homotopy}
Let $e_1,\ldots,e_n$ be exterior generators. The Koszul complex for $z_1^{d_1},\ldots,z_n^{d_n}$ is
\[
 C_q=R\otimes_\C\bigwedge^q\C^n,
 \qquad d_0=\sum_{i=1}^nz_i^{d_i}\iota_i,
\]
where $\iota_i$ is contraction by $e_i$, with the usual exterior signs. Regard $\PP$ as a complex concentrated in degree zero. Let $\iota:\PP\to C_0$ be inclusion, and set $p=0$ in positive degree.

Define $P_i$ on exterior monomials by applying $R_i$ to the coefficient if the monomial does not contain $e_i$, and by returning zero if it does. Define a degree-$+1$ operator
\begin{equation}\label{eq:h}
 h=\sum_{i=1}^n e_i\wedge D_iP_1\cdots P_{i-1}.
\end{equation}
An empty product is the identity. Every operation in this formula is defined on the original polydisk.

\begin{lemma}[Explicit contraction]\label{lem:contraction}
The operators satisfy
\begin{equation}\label{eq:contraction}
 d_0h+hd_0=1-\iota p,\qquad
 h^2=0,\quad h\iota=0,\quad ph=0,\quad p\iota=1.
\end{equation}
\end{lemma}
\begin{proof}
For the one-variable exterior complex, $h_i=e_i\wedge D_i$ satisfies
\[
 z_i^{d_i}\iota_i h_i+h_i z_i^{d_i}\iota_i=1-P_i:
\]
on degree zero this is $z_i^{d_i}D_i=1-R_i$, and on the component containing $e_i$ it is $D_i z_i^{d_i}=1$. Tensoring these identities in the ordered variables gives
\[
 d_0h+hd_0
 =\sum_{i=1}^n P_1\cdots P_{i-1}(1-P_i)
 =1-P_1\cdots P_n.
\]
Terms involving distinct exterior generators cancel with opposite signs. The final product is exactly $\iota p$.

For $h^2$, after an $e_i$ has been inserted, a later summand with larger index is killed by its preceding $P_i$, and a summand with the same index wedges $e_i$ twice. A summand with smaller index $j$ is killed by $D_jR_j=0$, because the first summand had already applied $P_j$. Thus $h^2=0$. A box monomial is killed by every relevant $D_i$, proving $h\iota=0$. The remaining relations follow directly from exterior degree and the definition of the box projection.
\end{proof}

\begin{example}[Two variables]
For degrees $(d_1,d_2)$, the degree-zero part is
\[
 h(f)=e_1D_1f+e_2D_2R_1f,
 \qquad f=R_1R_2f+z_1^{d_1}D_1f+z_2^{d_2}D_2R_1f.
\]
On $e_1g+e_2k$, one has $h(e_1g+e_2k)=e_1\wedge e_2D_1k$. The asymmetry records an order of division; it does not assert that the equations are triangular.
\end{example}

\section{Support-controlled deformation and finite freeness}\label{sec:division}
Extend all preceding operators coefficientwise to $\HH_\Gamma(V\times D)$ and its finite exterior powers. Put
\[
 F_i=z_i^{d_i}+E_i,\qquad
 T=\bigcup_i\supp(E_i)\subseteq\Gamma_{>0}.
\]
The union is well ordered because it is finite. The case $E_i=0$ is allowed. Let
\[
 \delta=\sum_iE_i\iota_i,\qquad d=d_0+\delta.
\]
Since multiplication by the $F_i$ commutes, $d^2=0$; more specifically $\delta^2=0$ and $d_0\delta+\delta d_0=0$.

\begin{theorem}[Perturbed contraction and canonical division]\label{thm:division}
Define
\begin{equation}\label{eq:perturbedops}
 Q=1+\delta h,\quad
 S=Q^{-1}=\sum_{k\ge0}(-\delta h)^k,\quad
 p_F=pS,\quad h_F=hS.
\end{equation}
These operators are well defined by strong summability and satisfy
\begin{equation}\label{eq:perturbedcontraction}
 p_Fd=0,\qquad p_F\iota=1,\qquad
 dh_F+h_Fd=1-\iota p_F.
\end{equation}
For every $H\in\HH_\Gamma(V\times D)$,
\begin{equation}\label{eq:division}
 H=\NF_F(H)+\sum_{i=1}^nF_iG_i,\qquad
 \NF_F(H)=p_FH\in\bigoplus_{\alpha\in\BB}\HH_\Gamma(V)z^\alpha,
\end{equation}
where $\sum_iG_ie_i=h_FH$. The remainder is the unique box polynomial congruent to $H$ modulo $(F_1,\ldots,F_n)$. The supports of the constructed remainder and quotients lie in $\supp(H)+T^*$.
\end{theorem}
\begin{proof}
Each occurrence of $\delta h$ introduces one multiplier from $T$, so \cref{lem:operator} constructs $S$ and gives the support bound. The key identity is an actual conjugation:
\begin{equation}\label{eq:conjugation}
 dQ=Qd_0.
\end{equation}
Indeed,
\[
 dQ=d_0+\delta+d_0\delta h+\delta^2h
     =d_0+\delta-\delta d_0h,
\]
whereas the original contraction and $\delta\iota=0$ give
\[
 Qd_0=d_0+\delta hd_0
      =d_0+\delta(1-\iota p-d_0h)
      =d_0+\delta-\delta d_0h.
\]
Thus $d=Qd_0S$ and $Sd=d_0S$. Also $Q\iota=\iota$ and $Qh=h$, by $h\iota=h^2=0$. Consequently
\[
 dh_F+h_Fd
 =Q(d_0h+hd_0)S
 =Q(1-\iota p)S
 =1-\iota pS.
\]
The other two relations in \eqref{eq:perturbedcontraction} follow from $pd_0=0$ and $S\iota=\iota$.

Apply the contraction to a degree-zero $H$, where $dH=0$, to obtain \eqref{eq:division}. If a box polynomial lies in the ideal, applying $p_F$ kills its ideal representation and fixes the polynomial. It is therefore zero. This proves uniqueness without presupposing the absence of syzygies. Finally all coefficientwise operations preserve the common ordinary domain, and the support estimate is the one already proved for $S$.
\end{proof}

\begin{corollary}[Finite free parameter algebra]\label{cor:free}
There is an isomorphism of $\HH_\Gamma(V)$-modules
\begin{equation}\label{eq:free}
 \mathcal A_F:=\frac{\HH_\Gamma(V\times D)}{(F_1,\ldots,F_n)}
 \ \cong\ \bigoplus_{\alpha\in\BB}\HH_\Gamma(V)z^\alpha.
\end{equation}
In particular $\mathcal A_F$ is finite free, hence flat, of rank $N$. The Koszul complex of $F_1,\ldots,F_n$ has zero homology in positive degrees, with the explicit contracting formula \eqref{eq:perturbedcontraction}.
\end{corollary}
\begin{proof}
Every operator in the construction is $\HH_\Gamma(V)$-linear: it differentiates and truncates only the $z$ variables. Finite convolution justifies commuting a Hahn parameter scalar with $S$. Existence and uniqueness of remainders prove \eqref{eq:free}. If $c$ is a positive-degree cycle, the contraction gives $c=d(h_Fc)$.
\end{proof}

\subsection{An explicit computational formula}
On degree zero, let
\begin{equation}\label{eq:L}
 L(H)=\sum_{i=1}^nE_iD_iR_1\cdots R_{i-1}H.
\end{equation}
Then
\begin{equation}\label{eq:NFseries}
 \NF_F(H)=R_1\cdots R_n\sum_{k\ge0}(-L)^kH.
\end{equation}
The formula computes the remainder directly. The full exterior calculation is what proves that this remainder annihilates the \emph{whole ideal}, not merely the particular expression produced by a division pass. Ignoring syzygies would leave a genuine gap in a several-variable argument.

\begin{corollary}[Finite coefficient recursion]\label{cor:recursion}
Write $L=\sum_{\tau\in T}t^\tau L_\tau$ and $Y=SH$. On the well-ordered set $\supp(H)+T^*$ the coefficients satisfy
\begin{equation}\label{eq:recursion}
 Y_\nu=H_\nu-\sum_{\substack{\tau\in T\,,\ \eta\in\supp(H)+T^*\\ \tau+\eta=\nu}}L_\tau Y_\eta,
 \qquad (\NF_FH)_\nu=pY_\nu.
\end{equation}
Each displayed sum is finite and every $\eta$ on the right is smaller than $\nu$.
\end{corollary}
\begin{proof}
Take the coefficient of $\nu$ in $(1+L)Y=H$ and apply the support lemma. Positivity of $\tau$ gives well-founded dependence. Equivalently, expand into the finitely many contributing words in \eqref{eq:NFseries}.
\end{proof}

This is a mathematical coefficient algorithm. An implementation additionally needs a representation of the support and access to the finitely many relevant decompositions. An arbitrary set-sized well-ordered support is not automatically an effective input format.

\subsection{Restrictions and specialization}
The operators commute with restriction of $V$, restriction to smaller centered polydisks, and enlargement of $\Gamma$. At an ordinary parameter $s=c$, normal forms specialize to those of the specialized system. At a finite parameter $s=c+\eta$, Taylor evaluation introduces only the positive monoid generated by $\supp\eta$, so the same statement holds after enlarging the workspace if necessary. This follows directly from finite contribution and commuting $z$-division with parameter derivatives.

No uniform bound on $E_i$ over $V$ is required. The assertion concerns common holomorphy of the coefficients, not a sup-norm-small analytic family on all of $V$.

\begin{corollary}[A global coherent implicit system]\label{cor:implicit}
If every $d_i=1$, there is a unique tuple $a(s)$ of positive-support Hahn-holomorphic functions on $V$ such that $F\sh(s,a(s))=0$ after evaluation. Every specialized infinitesimal zero is its specialization. Its support is contained in $T^*_{>0}$.
\end{corollary}
\begin{proof}
The quotient has basis $1$, so multiplication by $z_i$ is a scalar $a_i=\NF_F(z_i)$. Its exponent-zero coefficient vanishes, and its support is contained in $T^*$. The evaluation argument of \cref{sec:points} below proves that these scalars are exactly the specialized zeros. The identities can also be checked directly in \eqref{eq:NFseries}: substitution into an analytic coefficient function is justified by positive support and agrees with its action in the rank-one quotient.
\end{proof}

\section{From a finite quotient to actual surcomplex points}\label{sec:points}
For this section fix the parameters and write
\[
 A_F=\HH_\Gamma(D)/(F_1,\ldots,F_n).
\]
It is an $N$-dimensional $K$-algebra. The coefficientwise derivative of a coherent function evaluates to its fine derivative; equivalently one may work throughout with the strongly summable Taylor expansion at a point. In particular the local formal ideal at $a$ is generated by $F_i\sh(a+u)\in K[[u_1,\ldots,u_n]]$.

A finite-dimensional algebra alone is not yet a proof about actual analytic zeros. An arbitrary algebra homomorphism need not preserve an unspecified infinite summation operation. The next two lemmas close that gap explicitly.

\subsection{Finite jets at a fixed Hahn coefficient}
Let
\[
 b=\sum_{i=1}^n(d_i-1),\qquad d_{\max}=\max_i d_i.
\]
The total $z$-adic order of a holomorphic function means the least total degree of a nonzero Taylor monomial at zero.

\begin{lemma}[Finite-jet dependence]\label{lem:jets}
A term $pL^kH$ depends on the Taylor coefficients in $z$ of $H$ and of the perturbation coefficients involved only through total degree $b+kd_{\max}$. At a fixed Hahn exponent, every coefficient of $\NF_FH$, or of a multiplication matrix $M_H$, therefore depends on finitely many ordinary coefficient functions and finite jets of those functions.
\end{lemma}
\begin{proof}
An $R_i$ cannot lower total order, a $D_i$ lowers it by at most $d_i$, and multiplication by a holomorphic function cannot lower it. Thus $L^k$ lowers order by at most $kd_{\max}$. The projection $p$ kills every monomial of total degree greater than $b$. A Taylor tail initially of order greater than $b+kd_{\max}$ cannot survive. The same argument applies to a tail in one of the multiplier functions: there are at most $k$ subsequent division operators that could lower its order. This proves the jet bound for each fixed word. The support lemma permits only finitely many words at a fixed exponent. Multiplication matrices are obtained by applying normal form to $Hz^\beta$ for finitely many $\beta\in\BB$, so the conclusion applies to them as well.
\end{proof}

\subsection{A summable matrix functional calculus}
Let $X_i=M_{z_i}$ in the box basis. The matrices commute. Their entries belong to $K^\circ$, and their reductions modulo $\mm_K$ are the ordinary box-multiplication matrices $N_i$ for the quotient $\C[z]/(z_1^{d_1},\ldots,z_n^{d_n})$. In particular, any product of more than $b$ of the $N_i$ is zero.

\begin{lemma}[Matrix evaluation]\label{lem:matrix}
For $f\in\OO(D)$, with Taylor coefficients $f_\alpha$ at zero, the matrix family $\bigl(f_\alpha X^\alpha\bigr)_{\alpha\in\N^n}$ is strongly summable and
\begin{equation}\label{eq:matrixcalculus}
 M_f=\sum_{\alpha\in\N^n}f_\alpha X^\alpha.
\end{equation}
For $H=\sum_\gamma h_\gamma t^\gamma\in\HH_\Gamma(D)$ the corresponding doubly indexed matrix family is strongly summable, and its sum is $M_H$.
\end{lemma}
\begin{proof}
Write $X_i=N_i+Y_i$, with all entries of $Y_i$ of positive valuation. The positive supports of the finitely many $Y_i$ have a common well-ordered union $W$. Expand a matrix word into $N$ and $Y$ factors. A nonzero term with $k$ occurrences of $Y$ has at most $k+1$ intervening blocks of $N$ factors, each of length at most $b$. Its total length is therefore at most
\[
 (k+1)b+k.
\]
At a fixed exponent, the support lemma allows only finitely many $Y$ words and lengths. For each of these there are finitely many matrix words of the displayed bounded length. This proves strong summability; the support lies in $W^*$, or in $\supp(H)+W^*$ when the Hahn coefficients are included.

For a Taylor polynomial, \eqref{eq:matrixcalculus} is just the fact that multiplication in a quotient is a homomorphism. Fix a Hahn exponent. On the normal-form side, \cref{lem:jets} shows that sufficiently high Taylor terms contribute zero at that exponent. On the matrix-series side the preceding word argument gives the same eventual vanishing. The polynomial identity therefore proves the identity at that exponent. Repeat at each exponent, and use finite contribution to include the Hahn index $\gamma$.
\end{proof}

\begin{theorem}[Characters are exactly infinitesimal zeros]\label{thm:points}
The maximal ideals of $A_F$ are in bijection with the actual zeros of $F\sh$ in $D_K\sh$. Every zero lies in $\mm_K^n$. A character $\chi:A_F\to K$ corresponds to
\[
 a_i=\chi(z_i),\qquad \chi(H)=H\sh(a).
\]
In particular the zero set is nonempty and finite.
\end{theorem}
\begin{proof}
A maximal residue field of a finite-dimensional algebra over the algebraically closed field $K$ is $K$. For a character, each $a_i$ is a root of the characteristic polynomial of $X_i$. Reduction of $X_i$ is nilpotent, so this characteristic polynomial is monic of degree $N$ with all lower coefficients in $\mm_K$. Every root of such a polynomial is infinitesimal: a noninfinitesimal root $a$ would have finite inverse, and division of its equation by $a^N$ would give $1$ plus an infinitesimal equal to zero. Thus $a\in\mm_K^n$.

A character is a linear functional on the finite box basis. Such a functional commutes with strongly summable families of coefficient vectors: multiplication by each of its finitely many fixed scalar entries and finite addition preserve strong summability. Apply it to \cref{lem:matrix}, or equivalently apply that identity to the class of $1$ and then the character. This gives
\[
 \chi(H)=\sum_{\gamma,\alpha}t^\gamma
       \frac{\partial^\alpha h_\gamma(0)}{\alpha!}a^\alpha
       =H\sh(a).
\]
Since $F_i$ vanish in the quotient, the tuple $a$ is an actual zero. Conversely evaluation at a zero annihilates the ideal and gives a character. The box monomials generate the quotient, so the two assignments are inverse.

Finally, if $c+\varepsilon\in D_K\sh$ is a zero, taking standard parts gives $c_i^{d_i}=0$ for all $i$, hence $c=0$. The algebra is nonzero and finite-dimensional, so it has at least one and only finitely many maximal ideals.
\end{proof}

\subsection{Local multiplicity is preserved, not just the number of points}
For an actual zero $a$, define its formal analytic intersection algebra and multiplicity by
\begin{equation}\label{eq:localalg}
 Q_a=\frac{K[[u_1,\ldots,u_n]]}{
      (F_1\sh(a+u),\ldots,F_n\sh(a+u))},
 \qquad \mu_a=\dim_K Q_a.
\end{equation}
A priori one must prove that this dimension is finite and is the dimension of the corresponding local factor of $A_F$.

\begin{theorem}[Exact local intersection algebras]\label{thm:localalgebras}
Let $(A_F)_a$ be the Artinian local factor associated to the zero $a$. Taylor expansion at $a$ induces an isomorphism
\begin{equation}\label{eq:localiso}
 (A_F)_a\cong Q_a.
\end{equation}
Consequently
\begin{equation}\label{eq:conservation}
 \sum_{F\sh(a)=0}\mu_a=N.
\end{equation}
A zero is simple exactly when its Jacobian determinant is nonzero, and then $\mu_a=1$.
\end{theorem}
\begin{proof}
Let $C_i(T)=\det(TI-X_i)$. Cayley--Hamilton gives $C_i(z_i)=0$ in $A_F$, so canonical division gives an actual coherent ideal representation for $C_i(z_i)$. Taylor expansion of that identity at $a$ is legitimate by strong summability. Factor the ordinary polynomial over $K$ as
\[
 C_i(a_i+u_i)=u_i^{m_i}V_i(u_i),\qquad V_i(0)\ne0.
\]
The series $V_i$ is a formal unit. Hence $u_i^{m_i}=0$ in $Q_a$, and $Q_a$ is finite-dimensional and generated by polynomial classes. Taylor expansion therefore gives a surjection $A_F\to Q_a$.

In the local factor $(A_F)_a$, the elements $z_i-a_i$ are nilpotent. Every formal series in $u$ can consequently be evaluated there by a finite sum. The matrix functional calculus, regrouped at the infinitesimal point $a$, implies that $F_i\sh(a+u)$ evaluate to zero when $u_i$ is replaced by $z_i-a_i$. This regrouping is legitimate because the nilpotent variables allow only finitely many powers and the remaining powers of $a$ are strongly summable. Thus there is a homomorphism $Q_a\to(A_F)_a$.

The original surjection factors through $(A_F)_a$: a polynomial nonzero at $a$ maps to a unit in the local algebra $Q_a$, whereas the other Artinian factors are killed. The two maps are inverse on the coordinate generators. Those generators generate both algebras, so \eqref{eq:localiso} follows. The finite algebra decomposition
\[
 A_F\cong\prod_a(A_F)_a
\]
then gives \eqref{eq:conservation}. This decomposition follows, for example, by the Chinese remainder theorem applied to sufficiently high powers of the finitely many maximal ideals.

If the Jacobian is invertible, formal elimination by its invertible linear part gives $Q_a=K$. Conversely, if $Q_a=K$, its cotangent space is zero. The cotangent space has presentation $K^n$ modulo the span of the linear terms of the $F_i\sh(a+u)$, so the square Jacobian has full rank. This proves the simple-zero criterion.
\end{proof}

\begin{corollary}[A finite map of infinitesimal polydisks]\label{cor:finitemap}
The evaluated map
\[
 F\sh:\mm_K^n\longrightarrow\mm_K^n
\]
is surjective. Every fiber is finite and has total formal intersection multiplicity $N$. The same assertion holds in the full surcomplex infinitesimal class after allowing divisible workspace enlargement.
\end{corollary}
\begin{proof}
For any $c\in\mm_K^n$, the system $F_i-c_i$ has the same coordinate-monomial leading part and positive-support errors. Apply \cref{thm:points,thm:localalgebras}. The values of $F$ at infinitesimal points are infinitesimal by direct Taylor evaluation. For a full surcomplex target, choose a set-sized divisible workspace containing its supports and the original data. Conversely a zero in a larger workspace has coordinates satisfying the original characteristic polynomials; algebraic closedness of the smaller field shows that no additional zeros of a fixed system are introduced.
\end{proof}

Here ``finite map'' means the stated finite-fiber and finite-algebra properties. It does not assert properness for a topology on a proper class.

\section{Hahn residues and a perfect pairing}\label{sec:residues}
Choose ordinary $0<\rho_i<r_i$ and let
\[
 T_\rho=\{z\in\C^n:|z_i|=\rho_i\ \text{for all }i\}.
\]
It has product orientation $d\arg z_1\wedge\cdots\wedge d\arg z_n>0$. Define a coefficientwise contour functional by
\begin{equation}\label{eq:resdefinition}
 \RR_F(H)=\frac{1}{(2\pi i)^n}
 \HInt_{T_\rho}\frac{H(z)\dz}{F_1(z)\cdots F_n(z)}.
\end{equation}
Each inverse is expanded in positive Hahn support:
\begin{equation}\label{eq:denexpansion}
 \frac1{F_i(z)}=\sum_{k\ge0}\frac{(-E_i(z))^k}{z_i^{d_i(k+1)}}.
\end{equation}
The resulting coefficient functions are ordinary meromorphic functions near the torus, with possible poles only on coordinate hyperplanes. Integrate each coefficient using ordinary product contours. This is not integration along fine-continuous real-parameter paths.

\begin{theorem}[Residue expansion and ideal annihilation]\label{thm:residueformula}
The functional \eqref{eq:resdefinition} is well defined, independent of the radii, and $K$-linear. Explicitly,
\begin{equation}\label{eq:resseries}
 \RR_F(H)=\sum_{k\in\N^n}(-1)^{|k|}
 \bigl[z_1^{d_1(k_1+1)-1}\cdots z_n^{d_n(k_n+1)-1}\bigr]
 H\prod_{i=1}^nE_i^{k_i}.
\end{equation}
Its support is contained in $\supp(H)+T^*$. It annihilates $(F_1,\ldots,F_n)$ and thus induces a functional on $A_F$. With parameters, the same construction takes values in $\HH_\Gamma(V)$ and has the same properties over that ring.
\end{theorem}
\begin{proof}
At a fixed Hahn exponent, \cref{lem:support} makes the product of the expansions \eqref{eq:denexpansion} a finite sum of ordinary coefficient functions. Multivariable Cauchy coefficient extraction gives exactly \eqref{eq:resseries}, proving existence and independence of radii. Multiplication by a fixed Hahn scalar is legitimate by finite convolution, giving linearity and the support bound.

For a numerator $F_jH$, cancel $F_j$ before integrating. In every remaining coefficient, negative powers occur only in the variables $z_i$ with $i\ne j$. All factors are holomorphic in $z_j$ on its disk. The $z_j$ contour integral is therefore zero. Cancellation and rearrangement are justified at each exponent by the same support argument. Thus every ideal element is annihilated. Parameter holomorphy follows either by ordinary Cauchy differentiation under the compact contour or directly from Taylor coefficient extraction.
\end{proof}

\begin{theorem}[Perfect residue pairing]\label{thm:duality}
The pairing
\begin{equation}\label{eq:pairing}
 A_F\times A_F\longrightarrow K,\qquad
 (u,v)\longmapsto\RR_F(uv)
\end{equation}
is nondegenerate. In the parameter version it is a perfect pairing of finite free $\HH_\Gamma(V)$-modules. In particular
\begin{equation}\label{eq:membership}
 H\in(F_1,\ldots,F_n)
 \quad\Longleftrightarrow\quad
 \RR_F(Hz^\alpha)=0\quad\text{for all }\alpha\in\BB.
\end{equation}
\end{theorem}
\begin{proof}
In the box basis the Gram matrix is
\[
 G_{\alpha\beta}=\RR_F(z^{\alpha+\beta}).
\]
All entries have nonnegative Hahn support. At exponent zero, only $k=0$ in \eqref{eq:resseries} contributes, and
\[
 (G_0)_{\alpha\beta}=
 \begin{cases}1,&\alpha+\beta=(d_1-1,\ldots,d_n-1),\\0,&\text{otherwise}.
 \end{cases}
\]
This is a permutation matrix. Hence $\det G=\pm1+$ positive-support terms. It is a unit by a strongly summable geometric inverse, also over the parameter algebra because its leading coefficient is the nonzero constant $\pm1$. This proves perfection. For \eqref{eq:membership}, pass to the unique normal form. Its pairing against the basis vanishes exactly when its coefficient vector is zero.
\end{proof}

The residue is generally \emph{not} just the top coefficient of the normal form. That simplification holds for certain polynomial systems, including the example in \cref{sec:example}, but not for an arbitrary holomorphic error. The theorem uses the full, explicitly invertible Gram matrix.

\section{Finite-support transfer and the B\'ezoutian kernel}\label{sec:transfer}
The main division and multiplicity theorems above are proved directly. To transfer the normalization-sensitive classical residue identities, we need one additional argument. It reduces a \emph{single output coefficient} to a finite ordinary analytic deformation; it does not turn a Hahn field into a field with a real-valued norm.

\subsection{A finite-support transfer lemma}
\begin{lemma}[Transfer of coefficient identities]\label{lem:transfer}
Consider an expression obtained from positive-support errors by finitely many additions, products, derivatives, divided differences, box normal forms, finite matrices, and the residue expansion \eqref{eq:resseries}. Inverses of matrices with invertible ordinary constant leading term are also allowed. If an identity between such expressions holds for every sufficiently small finite-parameter ordinary holomorphic deformation of the coordinate-monomial system, then it holds for arbitrary well-ordered positive-support Hahn errors. Numerators may have arbitrary well-ordered Hahn support.
\end{lemma}
\begin{proof}
Fix one exponent $\nu$ of a proposed identity. Each constituent construction has a support bound obtained by adding input supports and positive generated monoids. Expanding the operator and geometric inverses gives only finitely many contributing words at $\nu$. Finite matrices add only finite products and sums. Thus only finitely many perturbation exponents, finitely many numerator exponents, and finitely many coefficient functions can affect this coefficient. For normal forms, \cref{lem:jets} further reduces the calculation to finite Taylor data. Residue terms are finite Taylor coefficient extractions as well.

Discard every perturbation coefficient not occurring in these finitely many contributions. Replace the remaining monomials $t^{\tau_1},\ldots,t^{\tau_q}$ by independent ordinary complex parameters $u_1,\ldots,u_q$. Their exponent relations are not assumed independent: they will be restored by the final substitution. For each fixed ordinary numerator coefficient, the hypothesis is an analytic identity for $u$ near zero. Its ordinary Taylor coefficients therefore agree. Substituting $u_j=t^{\tau_j}$ and collecting the finitely many terms contributing to $\nu$ proves the required coefficient equality. Finally add the finitely many relevant numerator coefficients with their original Hahn weights.

For completeness, the operator formulas really are the Taylor formulas of the finite ordinary deformations. Work on a smaller closed polydisk strictly inside $D$, in the Banach space of functions holomorphic in its interior and continuous on its closure. Each $R_i$ is bounded by Cauchy estimates and finite truncation. The removable quotient $D_i$ is bounded by applying the maximum principle on the $z_i$ boundary; for example a bound $(d_i+1)\rho_i^{-d_i}$ suffices. Consequently $\delta h$ has norm less than one for sufficiently small $u$, and its Neumann series converges in this ordinary Banach space. The contour denominators are likewise nonzero and have normally convergent geometric expansions there. These facts justify the ordinary Taylor identities used before substitution. No corresponding Banach convergence is claimed for the Hahn substitution.
\end{proof}

The passage through a smaller closed polydisk is only a device in this proof of an identity. The actual Hahn operators and their output coefficient functions remain defined on the original open $D$.

\subsection{Classical input and its role}
We use the ordinary local residue theorem for an isolated finite holomorphic intersection, with ordered equations and volume form. At a simple zero $a$, the residue is $H(a)/J_F(a)$; cluster residues add, and a change of generators $G=BF$ multiplies the numerator by $\det B$. These normalization facts and their algebraic counterparts are reviewed in \cite[\S0]{CDS}; \cite{CD} gives a broader introduction.

For a finite ordinary perturbation of $z_i^{d_i}$, choose a smaller polydisk and make the parameters small enough that $|E_i|<|z_i|^{d_i}$ on its $i$th boundary face. The product torus then represents the residue cycle of the enclosed cluster. This is the local product-domain version of the multidimensional residue theorem, equivalently obtained by deformation of the associated Leray cycles. Extra small constant perturbations make the cluster simple off the critical-value locus. Identities verified there extend to singular parameters because the fixed-contour expressions and the bounded normal-form operators are holomorphic in the parameters. Thus the cited classical input concerns only ordinary finite deformations; \cref{lem:transfer} supplies the arbitrary-support step.

\subsection{The diagonal B\'ezoutian}
Introduce a second variable tuple $w$. Choose the telescoping divided differences
\begin{equation}\label{eq:bezentries}
 B_{ij}(z,w)=
 \frac{F_i(w_1,\ldots,w_{j-1},z_j,z_{j+1},\ldots,z_n)
      -F_i(w_1,\ldots,w_j,z_{j+1},\ldots,z_n)}{z_j-w_j}.
\end{equation}
The quotient has a removable singularity on $z_j=w_j$. Put
\[
 \mathcal D_F(z,w)=\det(B_{ij}(z,w)).
\]
Then
\begin{equation}\label{eq:bezrelation}
 F_i(z)-F_i(w)=\sum_jB_{ij}(z,w)(z_j-w_j),
 \qquad \mathcal D_F(z,z)=J_F(z).
\end{equation}
The coefficients are holomorphic on $D\times D$, and their support lies in $\{0\}\cup T$ before taking the finite determinant. Applying box normal form in both tuples gives a well-defined class in $A_F\otimes_K A_F$.

\begin{theorem}[Residue reproducing kernel]\label{thm:bezout}
For every $H\in\HH_\Gamma(D)$,
\begin{equation}\label{eq:reproduce}
 \NF_{F,z}\!\left(\RR_{F,w}
       (H(w)\mathcal D_F(z,w))\right)=\NF_F(H)(z).
\end{equation}
Equivalently, if $b_1,\ldots,b_N$ is the box basis and $b_1^*,\ldots,b_N^*$ its residue-dual basis, then
\begin{equation}\label{eq:casimir}
 [\mathcal D_F(z,w)]=\sum_{j=1}^N b_j(z)b_j^*(w)
       \quad\text{in }A_F\otimes_K A_F.
\end{equation}
Both identities hold over the parameter algebra as well.
\end{theorem}
\begin{proof}
First consider a sufficiently small finite ordinary deformation with simple zeros. If $a\ne b$ are two zeros, \eqref{eq:bezrelation} gives a nonzero vector $a-b$ in the kernel of $B(a,b)$, so $\mathcal D_F(a,b)=0$. On the diagonal it equals $J_F(a)$. The ordinary residue theorem therefore gives
\[
 \RR_{F,w}(H(w)\mathcal D_F(a,w))=H(a).
\]
For a simple cluster, evaluation at its points identifies the finite quotient with the product of its residue fields. Hence the normal-form identity holds. Add independent small constant parameters to the equations if needed to reach simple clusters. Holomorphic dependence extends the identity to all small parameters, including zero. The transfer lemma now proves \eqref{eq:reproduce} for the Hahn system.

In a finite free module with a perfect pairing, the tensor representing the identity endomorphism is uniquely $\sum_j b_j\otimes b_j^*$. Formula \eqref{eq:reproduce} says exactly that $[\mathcal D_F]$ represents that identity. This proves \eqref{eq:casimir}. The argument is valid at each ordinary parameter point; uniqueness of holomorphic coefficient functions, or the same transfer argument with parameters left passive, proves the parameter version.
\end{proof}

This proof isolates the genuinely new summability issue from the classical B\'ezoutian identity. The latter is not being claimed as a new theorem of ordinary complex or polynomial analysis.

\section{Trace, local residues, and invariance}\label{sec:trace}
\begin{theorem}[Jacobian trace and characteristic polynomial]\label{thm:trace}
For $H\in\HH_\Gamma(D)$,
\begin{align}
 \Tr(M_H)&=\RR_F(HJ_F),\label{eq:trace}\\
 \det(TI-M_H)&=\prod_{F\sh(a)=0}(T-H\sh(a))^{\mu_a}.
 \label{eq:charpoly}
\end{align}
In particular,
\begin{equation}\label{eq:degree}
 \RR_F(J_F)=N,
\end{equation}
with no positive-support correction. The trace identity holds over $\HH_\Gamma(V)$ in the parameter version.
\end{theorem}
\begin{proof}
Let $b_j^*$ be residue-dual to $b_j$. The coefficient of $b_j$ in $Hb_j$ is $\RR_F(Hb_jb_j^*)$, so
\[
 \Tr(M_H)=\RR_F\!\left(H\sum_jb_jb_j^*\right).
\]
Apply the multiplication map $A_F\otimes A_F\to A_F$ to \eqref{eq:casimir}. By \eqref{eq:bezrelation}, $\sum_jb_jb_j^*=[J_F]$. This proves \eqref{eq:trace}, also over the parameter ring.

On the Artinian factor $(A_F)_a$, an element $H$ is its scalar value $H\sh(a)$ plus a nilpotent. Multiplication by that nilpotent is nilpotent; hence the characteristic polynomial of $M_H$ on this factor is $(T-H\sh(a))^{\mu_a}$. The product decomposition and \cref{thm:localalgebras} prove \eqref{eq:charpoly}. Take $H=1$ in the trace identity to obtain \eqref{eq:degree}.
\end{proof}

\subsection{Localization at actual surcomplex intersections}
Let $e_a\in A_F$ be the idempotent equal to $1$ in $(A_F)_a$ and $0$ in all other local factors. Define
\begin{equation}\label{eq:localized}
 \Res_{a,F}(H)=\RR_F(e_aH).
\end{equation}
The right side uses any box representative of $e_a$; it is independent of the representative by ideal annihilation. This definition works at singular intersections without writing a divergent quotient by the Jacobian.

\begin{theorem}[Moving-intersection residue theorem]\label{thm:localres}
There are finitely many localized residues and
\begin{equation}\label{eq:localres}
 \RR_F(H)=\sum_{F\sh(a)=0}\Res_{a,F}(H),\qquad
 \Res_{a,F}(HJ_F)=\mu_aH\sh(a).
\end{equation}
Each localized residue pairing on $(A_F)_a$ is nondegenerate. If $a$ is simple,
\begin{equation}\label{eq:simpleres}
 \Res_{a,F}(H)=\frac{H\sh(a)}{J_F\sh(a)}.
\end{equation}
Thus when all roots are simple the product-torus functional equals the sum of the actual simple-intersection residues, not merely a formal residue attached to their common standard part.
\end{theorem}
\begin{proof}
The idempotents sum to $1$, giving the first formula. Applying the trace identity to $e_aH$ restricts the trace to the local factor, where its value is $\mu_aH\sh(a)$, proving the second. Distinct local factors are orthogonal for the residue pairing because their product is zero. Perfection of the global pairing therefore implies perfection on each factor. If $a$ is simple, the factor is $K$ and $J_F$ acts by the nonzero scalar $J_F\sh(a)$. The second formula gives \eqref{eq:simpleres}.
\end{proof}

The definition \eqref{eq:localized} deliberately supplies its own normalization through the fixed contour and the B\'ezoutian. Via the local isomorphism \eqref{eq:localiso}, this is the residue-dual functional of the ordered formal complete intersection. The diagonal kernel \eqref{eq:casimir} is the usual algebraic characterization of that normalization; this agrees with the formal Grothendieck residue convention described in \cite[\S0]{CDS}. Neither an infinitesimal ordinary-radius contour nor a choice of a path between fine components is needed.

\subsection{Change of generators and infinitesimal coordinates}
\begin{theorem}[Transformation laws]\label{thm:transform}
Suppose $B(z)=I+C(z)$ is a coherent matrix with $v_H(C)>0$, and set $G=BF$. Then $G$ has the same coordinate-monomial leading system and
\begin{equation}\label{eq:genchange}
 \RR_G(H\det B)=\RR_F(H).
\end{equation}
Suppose $\Phi(z)=z+\eta(z)$, where all coordinates of $\eta$ have positive support. Coherent Taylor composition defines $F\circ\Phi$ on $D$ and
\begin{equation}\label{eq:coordchange}
 \RR_{F\circ\Phi}\bigl((H\circ\Phi)\det D\Phi\bigr)=\RR_F(H).
\end{equation}
The induced map $\Phi\sh:\mm_K^n\to\mm_K^n$ is bijective.
\end{theorem}
\begin{proof}
The matrix $B$ has a strongly summable inverse, and $G_i-z_i^{d_i}$ has positive support. For a finite ordinary deformation with simple intersections, $J_G(a)=\det B(a)J_F(a)$ at a common zero. The simple residue formula proves \eqref{eq:genchange}. Both contour expressions are holomorphic in small ordinary deformation parameters, so the identity extends across the singular locus and then to Hahn data by \cref{lem:transfer}.

For composition, expand every ordinary coefficient function of $F$ at the ordinary variable $z$ in powers of the positive-support displacement $\eta$. The support is controlled by the original support plus the positive monoid of $\eta$, and the coefficients remain holomorphic on $D$. For sufficiently small finite ordinary parameters on a smaller polydisk, $\Phi$ is a coordinate change near the cluster. The ordinary residue substitution rule, or the simple-root formula and the chain rule followed by parameter continuation, gives \eqref{eq:coordchange}. Transfer again proves the Hahn identity.

For any $w\in\mm_K^n$, the equations $z_i+\eta_i(z)-w_i=0$ have degrees $d_i=1$. By \cref{cor:finitemap} there is precisely one solution, of multiplicity one. This proves bijectivity on the whole infinitesimal polydisk, not just on an unspecified smaller fine neighborhood.
\end{proof}

These transformation laws show that the residue construction is not an artifact of a particular infinitesimal presentation of the same finite intersection. The ordinary leading coordinates are still part of the hypotheses.

\section{Stability of algebras and sharp stability of roots}\label{sec:stability}
Two different stability questions must be distinguished. Normal-form coefficients and cluster residues can be insensitive to a high-valuation perturbation, while individual roots are sensitive to an almost singular Jacobian. We prove estimates for both.

For a vector or matrix over $K$, let its valuation be the minimum valuation of its entries, with the zero object assigned $+\infty$. Then $v(AB)\ge v(A)+v(B)$. The analogous convention applies to matrices and vectors of coherent functions.

\subsection{The algebraic data do not lose valuation}
\begin{theorem}[Valuation stability of division and residues]\label{thm:algebrastability}
Let $F_i=z_i^{d_i}+E_i$ and $\widetilde F_i=z_i^{d_i}+\widetilde E_i$ have positive-support errors on the same domain. Suppose
\[
 v_H(\widetilde E-E)\ge\sigma>0.
\]
For every coherent $H$,
\begin{align}
 v_H\bigl(\NF_{\widetilde F}H-\NF_FH\bigr)&\ge v_H(H)+\sigma,\label{eq:NFstable}\\
 v\bigl(\RR_{\widetilde F}(H)-\RR_F(H)\bigr)&\ge v_H(H)+\sigma.
 \label{eq:resstable}
\end{align}
The multiplication matrices of a fixed $H$ in the common box basis obey the first bound as well. In particular the coordinate matrices, their characteristic-polynomial coefficients, the residue Gram matrix, and its inverse change only at valuation at least $\sigma$.
\end{theorem}
\begin{proof}
Set $S_F=(1+\delta_Fh)^{-1}$ and similarly for $\widetilde F$. The resolvent identity gives
\[
 S_{\widetilde F}-S_F
 =S_{\widetilde F}(\delta_F-\delta_{\widetilde F})hS_F.
\]
All $S$ operators preserve a lower valuation bound, while the middle perturbation difference increases it by at least $\sigma$. Applying $p$ proves \eqref{eq:NFstable}, and applying it to $Hz^\beta$ proves the matrix bound.

For residues, subtract the products of denominator inverses in \eqref{eq:resdefinition} by a finite telescoping identity. Each summand contains one factor $E_i-\widetilde E_i$, with all other geometric factors of nonnegative Hahn valuation. Coefficientwise integration cannot decrease valuation. This proves \eqref{eq:resstable}. Coordinate and Gram entries have nonnegative valuation, so finite determinant expansions preserve the same difference bound. The two Gram inverses have nonnegative entries, and their resolvent identity proves the final assertion.
\end{proof}

\subsection{A high-rank Hensel estimate for actual simple roots}
\begin{theorem}[Sharp conditioned root stability]\label{thm:rootstability}
Let $F_i=z_i^{d_i}+E_i$ be as above, and let $a\in\mm_K^n$ be a simple zero. Put
\begin{equation}\label{eq:kappa}
 A=DF\sh(a),\qquad \kappa=\max\{0,-v(A^{-1})\}.
\end{equation}
Let $\Delta$ be a coherent vector on $D$ with
\begin{equation}\label{eq:threshold}
 v_H(\Delta)\ge\sigma>2\kappa.
\end{equation}
Then $\widetilde F=F+\Delta$ has exactly one zero $a+h$ in
\[
 \mathcal B_a=\{a+u:v(u)>\kappa\}.
\]
It is simple and satisfies
\begin{align}
 v(h)&\ge\sigma-\kappa,\label{eq:shiftbound}\\
 v\bigl(h+A^{-1}\Delta\sh(a)\bigr)&\ge2\sigma-3\kappa.
 \label{eq:errorbound}
\end{align}
All statements remain valid in arbitrary divisible Hahn workspaces, without a rank-one or sequential-completeness assumption.
\end{theorem}
\begin{proof}
Every coefficient in the Taylor expansion of $F\sh(a+u)$ has nonnegative valuation, because the coherent coefficients of $F$ have nonnegative support and $a$ is infinitesimal. The coefficients of $\Delta\sh(a+u)$ have valuation at least $\sigma$. Write
\[
 F\sh(a+u)=Au+Q(u),\qquad \ord_u Q\ge2,
\]
and
\[
 \Delta\sh(a+u)=\Delta\sh(a)+B u+Q_\Delta(u),\qquad
 \ord_u Q_\Delta\ge2.
\]
Here $v(B)\ge\sigma$ and all coefficients of $Q_\Delta$ have that lower bound.

Set $r=\sigma-\kappa$, so $r>\kappa\ge0$, and substitute $u=t^r y$. The normalized equation is
\begin{equation}\label{eq:normalizedHensel}
 t^{-r}A^{-1}\widetilde F\sh(a+t^ry)
 =y+b+P(y)=0,\qquad b=t^{-r}A^{-1}\Delta\sh(a),
\end{equation}
where $v(b)\ge0$ and all nonzero coefficients of $P$ have strictly positive valuation. Indeed the extra linear term has valuation at least $\sigma-\kappa>0$, and a degree-$k\ge2$ term from $Q$ has valuation at least
\[
 (k-1)r-\kappa\ge r-\kappa=\sigma-2\kappa>0.
\]
Terms from $Q_\Delta$ have an additional valuation gain.

We must also justify that \eqref{eq:normalizedHensel} is a coherent equation, not merely a formal one. The Taylor coefficients at $a$ have a common nonnegative support contained in the original support plus the positive monoid generated by $\supp a$. Multiplication by the fixed matrix $A^{-1}$ and by powers $t^{(k-1)r}$ has a well-ordered support with finite contribution. At a fixed exponent, only finitely many degrees $k$ contribute. Thus each coefficient function of $y$ is a polynomial and the normalized equation is Hahn-coherent on any ordinary $y$ polydisk.

Translate $y=-\st(b)+v$. The new leading system is $v_i$, and all remaining coefficients have positive support. The degree-one case of \cref{thm:points,thm:localalgebras} gives a unique infinitesimal $v$, hence a finite $y$ solving \eqref{eq:normalizedHensel}. This constructs a zero with $v(h)\ge r$, proving \eqref{eq:shiftbound} and placing it in $\mathcal B_a$.

To prove uniqueness on the whole stated ball, not merely on the rescaled monad used in the construction, take $u,w$ with $v(u),v(w)>\kappa$. Taylor divided differences give
\begin{equation}\label{eq:injectiveball}
 \widetilde F\sh(a+u)-\widetilde F\sh(a+w)
 =(A+C(u,w))(u-w),
\end{equation}
where
\[
 v(C(u,w))\ge\min\{v(u),v(w),\sigma\}>\kappa.
\]
Every nonlinear divided-difference term contains at least one coordinate of $u$ or $w$, and the perturbation derivative has valuation at least $\sigma$. These expansions are strongly summable. Since $v(A^{-1}C)>0$, the matrix $I+A^{-1}C$ is invertible by its geometric series. Equation \eqref{eq:injectiveball} proves injectivity, hence uniqueness. The same estimate for the derivative at $a+h$ proves that the new zero is simple.

Finally the root equation gives
\[
 h+A^{-1}\Delta\sh(a)=-A^{-1}\bigl(Q(h)+Bh+Q_\Delta(h)\bigr).
\]
Using $v(h)\ge r$ yields the lower bound
\[
 \min\{2r-\kappa,\,\sigma+r-\kappa,\,\sigma+2r-\kappa\}
 =2\sigma-3\kappa,
\]
because $\kappa\ge0$. This proves \eqref{eq:errorbound}.
\end{proof}

This is a Hensel-type theorem with a support-based existence proof. The appearance of twice the inverse-Jacobian loss is the familiar conditioning mechanism, not a new phenomenon unique to surreals. What matters here is that the proof also covers non-Archimedean value groups where a countable sequence of increasingly accurate approximations need not approach every valuation scale.

\begin{corollary}[Stable matching of an entire simple cluster]\label{cor:matching}
Suppose all $N$ zeros of $F$ are simple. Let $\kappa_a$ be as in \eqref{eq:kappa}, and suppose $v_H(\Delta)\ge\sigma>2\max_a\kappa_a$. Then the zeros of $F$ and $F+\Delta$ have a unique matching characterized by
\[
 v(\widetilde a-a)>\kappa_a.
\]
All perturbed zeros are simple, and each matched pair satisfies \eqref{eq:shiftbound}--\eqref{eq:errorbound} with its own $\kappa_a$.
\end{corollary}
\begin{proof}
The balls $\mathcal B_a$ are disjoint. Otherwise, for two intersecting balls, the center of the ball with larger valuation threshold lies in the other ball, by the ultrametric inequality. Applying the injectivity argument \eqref{eq:injectiveball} with $\Delta=0$ would put two original zeros in a ball on which $F$ is injective. This is impossible. The theorem constructs one simple perturbed root in each of these $N$ balls. The total multiplicity of the perturbed system is $N$, so there can be no other roots.
\end{proof}

\subsection{All three bounds are attained in the quadratic model}
\begin{proposition}[Sharpness]\label{prop:sharpness}
The strict inequality $\sigma>2\kappa$ cannot be replaced by $\sigma\ge2\kappa$ in a uniform simple-root matching theorem. The valuation bounds \eqref{eq:shiftbound} and \eqref{eq:errorbound} are also sharp.
\end{proposition}
\begin{proof}
Let $\alpha>0$, $\tau=t^\alpha$, and $F(x)=x^2-\tau^2$. Its roots are $\pm\tau$. At $a=\tau$, the inverse derivative is $(2\tau)^{-1}$, so $\kappa=\alpha$. The perturbation $\Delta=\tau^2$ has valuation $2\kappa$ and changes $F$ to $x^2$. Two simple roots merge into one double root.

For $\Delta=t^\sigma$ with $\sigma>2\alpha$, the root near $\tau$ is
\[
 \tau+h=\tau\left(1-\frac{t^\sigma}{\tau^2}\right)^{1/2}
 =\tau-\frac{t^\sigma}{2\tau}
       -\frac{t^{2\sigma}}{8\tau^3}+\cdots.
\]
The binomial series is strongly summable because $\sigma-2\alpha>0$. Its displayed nonzero leading coefficients show
\[
 v(h)=\sigma-\alpha,\qquad
 v\left(h+\frac{t^\sigma}{2\tau}\right)=2\sigma-3\alpha.
\]
Thus both lower bounds are attained. The examples embed in any number of variables by adding equations $z_i=0$.
\end{proof}

\section{A coupled family: collision, discriminant, and infinite residues}\label{sec:example}
Let $\tau=t^\alpha$ with $\alpha>0$ and let $s$ be an ordinary holomorphic parameter. Consider
\begin{equation}\label{eq:example}
 F_1(x,y)=x^2-\tau y,\qquad
 F_2(s,x,y)=y^2-\tau sx.
\end{equation}
The leading system is $(x^2,y^2)$, so its coherent quotient is free of rank four over the parameter algebra, with ordered basis
\[
 (1,x,y,xy).
\]
This remains true at $s=0$, where the number of distinct roots drops.

\subsection{Multiplication matrices and the collision discriminant}
The relations $x^2=\tau y$ and $y^2=\tau sx$ give
\begin{equation}\label{eq:examplematrices}
 X=\begin{pmatrix}
 0&0&0&0\\1&0&0&\tau^2s\\0&\tau&0&0\\0&0&1&0
 \end{pmatrix},\qquad
 Y=\begin{pmatrix}
 0&0&0&0\\0&0&\tau s&0\\1&0&0&\tau^2s\\0&1&0&0
 \end{pmatrix}.
\end{equation}
They commute and satisfy $X^2=\tau Y$, $Y^2=\tau sX$. Their first characteristic polynomial is
\begin{equation}\label{eq:examplechar}
 \det(TI-X)=T^4-\tau^3sT.
\end{equation}
The Jacobian is $J_F=4xy-\tau^2s$, and
\begin{equation}\label{eq:exampledisc}
 \det M_{J_F}=\det(4XY-\tau^2sI)=-27\tau^8s^4.
\end{equation}
Thus $s\ne0$ gives a reduced four-point fiber, whereas $s=0$ gives a nonreduced fiber. More generally, for any finite system in this article, $\det M_{J_F}\ne0$ is equivalent to all roots being simple: by \eqref{eq:charpoly}, it is the product of the Jacobian values with multiplicities, and \cref{thm:localalgebras} supplies the Jacobian criterion.

At $s=0$, elimination gives $y=x^2/\tau$ and $x^4=0$. Consequently the single root $(0,0)$ has local algebra $K[x]/(x^4)$ and multiplicity four. There is no loss of intersection length at the collision. Formula \eqref{eq:exampledisc} is a coherent function of the parameter that detects it.

\subsection{The residue pairing is explicit}
In this example,
\begin{equation}\label{eq:examplenormalres}
 \RR_F(H)=\text{the coefficient of }xy\text{ in }\NF_F(H).
\end{equation}
To check this, it is enough to evaluate residues on the box basis. Expanding the denominators gives
\[
 \frac1{(x^2-\tau y)(y^2-\tau sx)}
 =\sum_{k,l\ge0}\tau^{k+l}s^l
          x^{l-2k-2}y^{k-2l-2}.
\]
For a box numerator $x^py^q$, a residue term requires
\[
 p+l=2k+1,\qquad q+k=2l+1.
\]
Adding the equations gives $k+l=p+q-2$. Since $p,q\in\{0,1\}$, the only possible case is $p=q=1$, $k=l=0$. Thus the basis residues are $(0,0,0,1)$, proving \eqref{eq:examplenormalres} by ideal annihilation.

The residue Gram matrix is
\begin{equation}\label{eq:examplegram}
 G=\begin{pmatrix}
 0&0&0&1\\0&0&1&0\\0&1&0&0\\1&0&0&\tau^2s
 \end{pmatrix},\qquad \det G=1.
\end{equation}
In particular it stays perfect at the singular parameter $s=0$. Duality survives even when a formula using $1/J_F(a)$ at every root ceases to make sense.

\subsection{Four actual infinitesimal roots and a finite sum of infinite residues}
Now set $s=1$. Besides $(0,0)$, the roots are
\begin{equation}\label{eq:exampleroots}
 (x,y)=(\tau\zeta,\tau\zeta^2),\qquad \zeta^3=1.
\end{equation}
They are all simple. The Jacobian equals $-\tau^2$ at the origin and $3\tau^2$ at each nonzero root. Thus the residues of the constant numerator $1$ are
\[
 -\tau^{-2},\qquad \frac1{3}\tau^{-2},\qquad
 \frac1{3}\tau^{-2},\qquad \frac1{3}\tau^{-2}.
\]
Every residue is infinite, but their sum is exactly zero. For the numerator $xy$ their sum is $1$, consistent with \eqref{eq:examplenormalres}.

For a genuinely nonpolynomial numerator $H=e^{x+y}$, interpreted by ordinary holomorphic coefficients and infinitesimal Taylor evaluation, the moving-root formula yields
\begin{equation}\label{eq:exponentialres}
 \RR_F(e^{x+y})
 =\frac{e^{2\tau}+2e^{-\tau}-3}{3\tau^2}
 =\sum_{m\ge0}\frac{2^{m+2}+2(-1)^{m+2}}{3(m+2)!}\tau^m.
\end{equation}
Here $\zeta+\zeta^2=-1$ for the two primitive cube roots. The first terms are
\[
 \RR_F(e^{x+y})
 =1+\frac\tau3+\frac{\tau^2}{4}+\frac{\tau^3}{12}
   +\frac{11\tau^4}{360}+\frac{\tau^5}{120}+\cdots.
\]
The direct coefficientwise contour calculation gives a different-looking expression:
\begin{equation}\label{eq:doubleres}
 \RR_F(e^{x+y})
 =\sum_{\substack{k,l\ge0\\2k+1-l\ge0\\2l+1-k\ge0}}
  \frac{\tau^{k+l}}{(2k+1-l)!(2l+1-k)!}.
\end{equation}
The theorem proves equality to all orders. As a separate consistency check, the accompanying script compares all coefficients through $\tau^{12}$ exactly.

At all four roots, the least valuation of the inverse Jacobian entries is $-\alpha$, so $\kappa=\alpha$. Any coherent perturbation of valuation strictly larger than $2\alpha$ therefore matches all four roots uniquely, with displacement valuation at least $\sigma-\alpha$. At the same time, \cref{thm:algebrastability} shows that the full cluster residue of a numerator of valuation zero changes only at valuation $\sigma$ or higher. Individual roots lose $\alpha$ of accuracy; the cluster residue does not.

\section{Examples genuinely beyond a finite-parameter polynomial model}\label{sec:infinite}
The preceding family permits elementary elimination. The theorem also covers inputs for which no finite list of ordinary deformation parameters describes the whole function.

\begin{example}[Infinitely many scales below one valuation threshold]\label{ex:infinitesupport}
In a workspace containing $\Q$, let
\[
 \lambda_m=1-\frac1{m+1},\qquad m=1,2,\ldots,
\]
and define on an ordinary bidisk
\begin{align*}
 F_1(x,y)&=x^3+
      \sum_{m\ge1}t^{\lambda_m}e^{m!x}(1+y^m),\\
 F_2(x,y)&=y^2+
      \sum_{m\ge1}t^{\lambda_m}e^{m!y}(x+x^my).
\end{align*}
The support $\{\lambda_m:m\ge1\}$ is strictly increasing and well ordered, although infinitely many exponents lie below $1$. The coefficient functions are holomorphic with no uniform growth bound in $m$. Nevertheless the coherent quotient has basis
\[
 1,x,x^2,y,xy,x^2y
\]
and rank six. All six intersections, counted with local multiplicity, lie in the infinitesimal bidisk. Its residue pairing is perfect, and its multiplication and residue coefficients are computed by the explicit support-controlled formulas.
\end{example}

An algorithm that asks for ``all coefficients of valuation less than $1$'' has infinitely much output in this example. This does not contradict finite contribution at each \emph{specified} exponent. The two notions of truncation must not be conflated.

\begin{example}[A value group of higher rank]
Take the divisible group $\Gamma=\Q+\Q\omega\subseteq\No$. In the Hahn convention $t^\omega$ is smaller than every $t^m$, $m\in\N$. The scalar series
\[
 (1-t)^{-1}=\sum_{m\ge0}t^m
\]
is strongly summable, but the sequence of finite partial sums does not approach its sum to valuation accuracy $\omega$. In a coupled system containing both $t$ and $t^\omega$ perturbations, a countable geometric iteration cannot be justified by saying that its errors reach every scale. The contraction and residue constructions still work, because a fixed exponent has finitely many contributing positive words.
\end{example}

\subsection{Arbitrary ordinary degrees in the perturbation}
The errors need not be lower-degree polynomials in $z$. For example,
\[
 F_1=x^2+t\sin(e^{x+y}),\qquad F_2=y^3+t^{\sqrt2}e^{xy}
\]
have a rank-six quotient on every ordinary bidisk on which the coefficients are considered. The ordinary entire expressions in this display do not imply a chosen global surcomplex sine or exponential at arbitrary infinite arguments. Only their ordinary holomorphic coefficient functions and their infinitesimal Taylor evaluations are used. The degree of the reduction controls the finite monad intersection; there is no assertion about zeros at infinite points outside that domain.

\section{Computation, verification, and proof dependencies}\label{sec:verification}
\subsection{What the formulas compute}
To compute a selected coefficient of a normal form, expand \eqref{eq:NFseries} into the finitely many positive-support words that can reach that exponent. Apply only the finite Taylor jets required by \cref{lem:jets}, and then apply the box projection. The quotients in the division certificate are obtained from $h_FH$. Multiplication matrices require this calculation only for the finitely many inputs $Hz^\beta$, $\beta\in\BB$.

Cluster residues can be computed by \eqref{eq:resseries}, without solving for any roots. Local residues require the Artinian factorization or the corresponding idempotents. When all roots are simple, the formula $H(a)/J_F(a)$ is equivalent but can expose large canceling quantities. The matrix trace formula offers a third exact representation. Which is computationally preferable depends on the input support and on the complexity of ordinary coefficient extraction.

No general complexity bound is claimed for arbitrary noncomputable supports or arbitrary holomorphic-function oracles. In a finitely generated positive monoid, or for an explicitly enumerable support with decidable finite decompositions, the formulas become concrete symbolic procedures. The high-rank examples show why a naive integer-degree loop is not a general implementation.

\subsection{Reproducible symbolic checks}
The archive contains \texttt{verify\_examples.py} and \texttt{verification\_report.txt}. With Python and SymPy, the script performs 58 exact checks:
\begin{enumerate}
\item The monomial Koszul contraction, $h^2=0$, and $dQ=Qd_0$ in all eight exterior components of a three-variable example.
\item A coupled division certificate, ideal annihilation on the box basis, and commuting multiplication matrices, modulo $\tau^6$.
\item The two-variable example's exact matrices, defining relations, characteristic polynomial, Jacobian discriminant, perfect residue pairing, and trace-residue formula on a basis.
\item Equality of the two exponential-residue calculations through $\tau^{12}$.
\end{enumerate}
All checks passed in the recorded run. The script uses the symbol \texttt{a} for the ordinary parameter denoted $s$ in the article. It checks specified finite identities and truncations, not the abstract theorems for all supports.

\subsection{A support and dependency audit}
\begin{center}\small
\begin{longtable}{@{}L{0.30\textwidth}L{0.59\textwidth}@{}}
\toprule
Construction & Controlling fact\\
\midrule
\endhead
Infinitesimal evaluation & $S+(\bigcup_i\supp\varepsilon_i)^*$; finite multivariate Taylor multiplicity.\\[0.3em]
Perturbed inverse and division & $S+T^*$; each occurrence of $\delta h$ introduces a positive letter.\\[0.3em]
Quotient freeness & Explicit conjugation $dQ=Qd_0$, not an assumed independence of remainders.\\[0.3em]
Analytic matrix evaluation & Finitely many positive matrix factors at each exponent; bounded intervening nilpotent blocks.\\[0.3em]
Characters and actual roots & Matrix Taylor calculus; no implicit assumption that a character preserves an arbitrary infinite sum.\\[0.3em]
Local intersection multiplicity & Cayley--Hamilton ideal identities and inverse maps between the Artinian factor and the formal local quotient.\\[0.3em]
Residue functional & $S+T^*$; geometric denominator expansion and finite Cauchy coefficient extraction.\\[0.3em]
Residue duality & Gram matrix has invertible ordinary anti-diagonal reduction.\\[0.3em]
B\'ezoutian and transformation identities & Classical finite-parameter identities, transferred one Hahn coefficient at a time.\\[0.3em]
Root stability & Rescaling makes the normalized nonlinear part strictly positive; degree-one finite-map theorem supplies existence.\\
\bottomrule
\end{longtable}
\end{center}

The most delicate bridges are the matrix functional calculus and the finite-support transfer lemma. The former is needed because finite dimension by itself does not identify abstract characters with analytic evaluation. The latter is needed because ordinary continuity in a real parameter does not automatically justify a Hahn substitution. Neither bridge is replaced by a numerical experiment.

\section{Scope, extensions, and remaining mathematical questions}\label{sec:scope}
\subsection{A controlled-resolution extension criterion}
The monomial leading equations were used to write an explicit fixed-domain contraction. The deformation argument itself is more general.

\begin{proposition}[Abstract controlled-resolution criterion]\label{prop:abstract}
Let an ordinary Koszul complex of holomorphic functions $f_1,\ldots,f_n$ on a fixed domain admit complex-linear maps $h,p,\iota$ satisfying
\[
 d_0h+hd_0=1-\iota p,\quad h^2=h\iota=p h=0,\quad p\iota=1,
\]
where the target of $p$ is a finite-dimensional space in degree zero, and all maps preserve the chosen ordinary holomorphic function spaces. Then every positive-support Hahn perturbation $f_i+E_i$ has a finite-dimensional quotient of the same dimension, with a canonical remainder and the same $S+T^*$ support bound. If parameter-linearity is part of the original contraction, the corresponding relative finite-free statement holds.
\end{proposition}
\begin{proof}
Extend the maps coefficientwise. Multiplication by each $E_i$ introduces a positive letter, so the operator inverse exists by \cref{lem:operator}. The proof of $dQ=Qd_0$ uses only the displayed contraction identities, anticommutation of the Koszul differentials, and $\delta\iota=0$. Therefore the proof of \cref{thm:division,cor:free} applies verbatim.
\end{proof}

This criterion is an algebraic extension, not a claim that an arbitrary choice of a vector-space splitting is analytically satisfactory. To recover the actual-point and residue theorems for a new leading singularity, one also needs a finite-jet or equivalent continuity property sufficient for \cref{lem:matrix,lem:transfer}, and a compatible residue normalization. Our coordinate-monomial contraction verifies these properties explicitly.

\subsection{What has been resolved within the stated category}
For coordinate-monomial leading complete intersections, the several-variable division and finite-mapping target of the supplied manuscript is answered affirmatively. The result includes parameter-dependent finite freeness, exact local intersection multiplicities, a residue-dual algebra, coordinate and generator compatibility for positive-support changes, and sharp stability of simple fibers. These conclusions apply to infinite Hahn support and to value groups of arbitrary rank. They are not restricted to perturbations that happen to be a convergent one-parameter complex family.

\subsection{Questions left genuinely open by this article}
A natural next problem is to construct explicit fixed-domain, finite-jet-controlled Koszul contractions for general isolated ordinary complete intersections, with useful parameter-uniform versions. Ordinary local analytic division suggests a route, but its domain and support consequences must be proved rather than assumed. The abstract criterion specifies exactly what additional input would extend the present result.

A second problem is effective high-rank computation: given a structured support presentation, determine finite-decomposition algorithms and quantitative jet bounds that are practical for multiplication matrices and discriminants. Infinite support below one valuation cut prevents the usual finite-precision model from applying unchanged.

A third problem is global geometry. The finite monad maps here provide local models, but a coherent atlas across ordinary domains and changing surreal scales requires compatibility of supports, residue classes, and branch data. The all-scale rigidity phenomenon in \cite{Supplied} shows why a naive atlas condition can exclude transcendental functions altogether. This article does not choose a universal extension mechanism across essential singularities.

\subsection{Priority and mathematical confidence}
The priority check examined retrieved primary-source papers and publication records, together with targeted searches for surcomplex finite mappings, Hahn-coherent division, and residue/intersection analogues. No source located in that search stated the combined fixed-domain, arbitrary-support theorem in the precise category used here. That is evidence for investigating priority, not a proof of priority. The classical antecedents are close and important, especially homological perturbation and the residue--normal-form correspondence; terminology alone would not make a result new.

The proofs in this manuscript are intended to stand independently of a novelty claim. The monomial reduction hypothesis, the strong-summability convention, and the classical finite-parameter residue input are explicit. Independent review would be particularly valuable for the analytic-to-algebraic bridges identified in \cref{sec:verification}. No assertion of peer review or formal proof checking is made.

\section{Conclusion}
A finite intersection in surcomplex analysis can be more than a count of roots with a common standard part. Under an explicit leading-system hypothesis, it has a finite free analytic algebra, an exact local multiplicity decomposition, and a residue pairing that survives collisions. The contraction
\[
 h_F=h(1+\delta h)^{-1}
\]
constructs that algebra on the original ordinary domain, while finite support at each exponent replaces unavailable sequential convergence. The B\'ezoutian then links the contour functional to multiplication matrices and to the actual infinitesimal intersections.

The stability results separate robust cluster data from sensitive individual roots. Algebraic coefficients and residues preserve perturbation valuation; root locations lose precisely the amount measured by the inverse Jacobian. A quadratic model attains the resulting bounds and the collision threshold. Together these results supply a concrete several-variable continuation of the one-variable coherent theory, with substantial room for extension beyond coordinate-monomial leading singularities.

\appendix
\section{Notation and conventions}
\begin{center}\small
\begin{tabular}{@{}L{0.25\textwidth}L{0.64\textwidth}@{}}
\toprule
Notation & Meaning\\
\midrule
$\SC=\No[i]$ & Complexification of the surreal ordered class field.\\[0.25em]
$K=\C((t^\Gamma))$ & Set-sized divisible Hahn workspace inside $\SC$.\\[0.25em]
$t^\gamma$ & Conway normal-form monomial $\omega^{-\gamma}$.\\[0.25em]
$v$, $v_H$ & Least exponent of a scalar or a nonzero coefficient function.\\[0.25em]
$K^\circ$, $\mm_K$ & Finite elements and infinitesimals in the workspace.\\[0.25em]
$\HH_\Gamma(U)$ & Common well-ordered Hahn families in $\OO(U)$.\\[0.25em]
$H\sh$ & Evaluation by strongly summable infinitesimal Taylor expansion.\\[0.25em]
$\BB$ & Multi-indices $0\le\alpha_i<d_i$; $|\BB|=N$.\\[0.25em]
$R_i,D_i$ & Taylor truncation and division by $z_i^{d_i}$.\\[0.25em]
$d_0,\delta,d$ & Initial, perturbing, and full Koszul differentials.\\[0.25em]
$h,S,h_F,p_F$ & Initial contraction, inverse $S=(1+\delta h)^{-1}$, and perturbed operators.\\[0.25em]
$A_F$, $Q_a$, $\mu_a$ & Global cluster algebra, formal local algebra, and its dimension.\\[0.25em]
$X_i,M_H$ & Multiplication by a coordinate and by an arbitrary numerator.\\[0.25em]
$\RR_F$, $\Res_{a,F}$ & Product-torus Hahn residue and its Artinian localization.\\[0.25em]
$\kappa$ & Inverse-Jacobian valuation loss at a simple root.\\
\bottomrule
\end{tabular}
\end{center}
All derivatives are with respect to the analytic variables. They do not differentiate the Hahn monomials as transseries. An inequality between valuations is an inequality in the ordered group $\Gamma$, not an estimate in an unstated real norm.

\begin{thebibliography}{99}
\addcontentsline{toc}{section}{References}
\bibitem{Supplied}
\emph{Surcomplex Analysis: Infinitesimal Calculus, Hahn-Coherent Holomorphy, and Contour Theory}.
Unpublished manuscript supplied with this research request, dated September 21, 2026;
source file \texttt{surcomplex\_analysis(2).tex}.
Used as the starting framework and as the source of the several-variable research target; authorship and publication status are not inferred.

\bibitem{Alling}
Norman L. Alling,
\emph{Foundations of Analysis over Surreal Number Fields},
North-Holland Mathematics Studies 141, North-Holland, 1987.

\bibitem{BM}
Alessandro Berarducci and Vincenzo Mantova,
``Transseries as germs of surreal functions,''
\emph{Transactions of the American Mathematical Society} \textbf{371} (2019), 3549--3592.
\href{https://arxiv.org/abs/1703.01995}{arXiv:1703.01995}.

\bibitem{CD}
Eduardo Cattani and Alicia Dickenstein,
``Introduction to residues and resultants,''
in A. Dickenstein and I. Z. Emiris (eds.),
\emph{Solving Polynomial Equations: Foundations, Algorithms, and Applications},
Algorithms and Computation in Mathematics 14, Springer, 2005, pp.~1--61.
\href{https://doi.org/10.1007/3-540-27357-3_1}{doi:10.1007/3-540-27357-3\_1}.

\bibitem{CDS}
Eduardo Cattani, Alicia Dickenstein, and Bernd Sturmfels,
``Computing multidimensional residues,''
in \emph{Algorithms in Algebraic Geometry and Applications},
Progress in Mathematics 143, Birkh\"auser, 1996, pp.~135--164.
\href{https://arxiv.org/abs/alg-geom/9404011}{arXiv:alg-geom/9404011}.
The numbering cited in the text is that of the arXiv manuscript, especially \S0.

\bibitem{CE}
Ovidiu Costin and Philip Ehrlich,
``Integration on the surreals,''
\emph{Advances in Mathematics} \textbf{452} (2024), article 109823.
\href{https://arxiv.org/abs/2208.14331}{arXiv:2208.14331}.

\bibitem{Cox}
David A. Cox,
\emph{Solving Equations via Algebras},
CIMPA Graduate School lecture notes, Buenos Aires, 2003.
\href{https://mate.dm.uba.ar/~visita16/cimpa/notes/Cox.pdf}{Author's course notes}.
An expanded chapter appears in \emph{Solving Polynomial Equations}, Springer, 2005.

\bibitem{Crainic}
Marius Crainic,
\emph{On the Perturbation Lemma, and Deformations}, 2004.
\href{https://arxiv.org/abs/math/0403266}{arXiv:math/0403266}.

\bibitem{MS}
J\"orn M\"uller and Alexander Strohmaier,
``The theory of Hahn meromorphic functions, a holomorphic Fredholm theorem and its applications,''
\emph{Analysis \& PDE} \textbf{7} (2014), 745--770.
\href{https://arxiv.org/abs/1205.0236}{arXiv:1205.0236};
\href{https://msp.org/apde/2014/7-3/apde-v7-n3-p08-s.pdf}{publisher full text}.

\bibitem{Neumann}
B. H. Neumann,
``On ordered division rings,''
\emph{Transactions of the American Mathematical Society} \textbf{66} (1949), 202--252.

\bibitem{Poonen}
Bjorn Poonen,
``Maximally complete fields,''
\emph{L'Enseignement Math\'ematique} (2) \textbf{39} (1993), 87--106.
\href{https://math.mit.edu/~poonen/papers/amsval.pdf}{Author's full text}.

\bibitem{vdDE}
Lou van den Dries and Philip Ehrlich,
``Fields of surreal numbers and exponentiation,''
\emph{Fundamenta Mathematicae} \textbf{167} (2001), 173--188;
erratum \textbf{168} (2001), 295--297.
\href{https://doi.org/10.4064/fm167-2-3}{doi:10.4064/fm167-2-3}.
\end{thebibliography}
\end{document}
